\documentclass[final]{beamer}
\usepackage[size=a0,scale=1.14]{beamerposter}
\usetheme{default}
\usecolortheme{dove}
\usefonttheme{professionalfonts}

\usepackage[T1]{fontenc}
\usepackage[utf8]{inputenc}
\usepackage{lmodern}
\usepackage{amsmath,amssymb,mathtools,bm}
\usepackage{booktabs}
\usepackage{graphicx}
\usepackage{tikz}
\usetikzlibrary{arrows.meta,positioning,calc}
\usepackage{xcolor}
\usepackage{hyperref}

\definecolor{imaBlue}{RGB}{0,56,101}
\definecolor{imaTeal}{RGB}{0,124,146}
\definecolor{imaGray}{RGB}{245,247,250}
\definecolor{imaDark}{RGB}{40,44,52}

\setbeamercolor{background canvas}{bg=white}
\setbeamercolor{block title}{fg=white,bg=imaBlue}
\setbeamercolor{block body}{fg=imaDark,bg=imaGray}
\setbeamercolor{block title alerted}{fg=white,bg=imaTeal}
\setbeamercolor{block body alerted}{fg=imaDark,bg=white}
\setbeamercolor{headline}{fg=white,bg=imaBlue}
\setbeamercolor{footline}{fg=white,bg=imaBlue}

\setbeamertemplate{navigation symbols}{}
\setbeamertemplate{headline}{
  \begin{beamercolorbox}[wd=\paperwidth,ht=5.3ex,dp=1.8ex,leftskip=1.3cm,rightskip=1.3cm]{headline}
    {\LARGE \textbf{Motion-Enabled Tomography via Gaussian Mixture Models}}\hfill
    {\large IMA Conference on Inverse Problems: From Theory to Application}
  \end{beamercolorbox}
  \begin{beamercolorbox}[wd=\paperwidth,ht=4.2ex,dp=1.2ex,leftskip=1.3cm,rightskip=1.3cm]{headline}
    {\large Daniel Burrows$^{1}$ \quad Can Evren Yarman$^{2}$ \quad Ozan \"Oktem$^{3}$}
    \hfill
    {\large $^{1}$University of Bath \quad $^{2}$SLB \quad $^{3}$KTH Royal Institute of Technology}
  \end{beamercolorbox}
}

\setbeamertemplate{footline}{
  \begin{beamercolorbox}[wd=\paperwidth,ht=3.6ex,dp=1.4ex,leftskip=1.3cm,rightskip=1.3cm]{footline}
    \href{https://github.com/dwb26/gmm-ct}{github.com/dwb26/gmm-ct}
    \hfill
    Contact: dwb26@bath.ac.uk
  \end{beamercolorbox}
}

\newcommand{\R}{\mathbb{R}}
\newcommand{\RayTran}{\mathcal{X}}
\newcommand{\argmin}{\operatorname*{arg\,min}}
\newcommand{\argmax}{\operatorname*{arg\,max}}
\newcommand{\I}{\mathbb{I}}
\newcommand{\resdir}{data/results/20260526_145626_seed9_N5}

\begin{document}
\begin{frame}[t]
\begin{columns}[t,totalwidth=\textwidth]

\begin{column}{0.32\textwidth}

\begin{block}{Problem and Motivation}
\textbf{Goal.} Reconstruct moving objects in CT using a fixed source-detector geometry.

\vspace{0.4em}
\textbf{Challenge.} The inverse problem is severely ill-posed when motion is unknown.

\vspace{0.4em}
\textbf{Approach.} Jointly recover:
\begin{itemize}
\item projectile motion,
\item rotational motion,
\item object morphology,
\end{itemize}
using a Gaussian mixture model (GMM) representation.

\vspace{0.4em}
\textbf{Why GMM?}
\begin{itemize}
\item expressive approximation in object space,
\item closed-form ray transform,
\item exact gradients for optimization.
\end{itemize}
\end{block}

\begin{block}{Spatiotemporal Forward Model}
Represent the object as
\[
 f(\bm{x},t)=\sum_{n=1}^{N}\alpha_n\bigl(M_n(t).U_n^{-1}.\rho_0\bigr)(\bm{x}),
\]
with isotropic template $\rho_0(\bm{x})=\exp(-\bm{x}^\top\bm{x})$.

Data are ray measurements
\[
 g(\ell,t)=\RayTran\bigl(f(\cdot,t)\bigr)(\ell),\qquad
 \RayTran(f)[\bm{s},\bm{r}] = \int_{\R} f\bigl(\bm{s}+\xi\,\widehat{\bm{r}-\bm{s}}\bigr)\,d\xi.
\]

For anisotropic Gaussian $A=(U^{-1},\mu)$:
\[
\RayTran(A.\rho_0)[\bm{s},\bm{r}] =
\frac{\sqrt{\pi}}{\|U\widehat{\bm{r}-\bm{s}}\|}
\exp\!\left(
\frac{\langle U(\bm{r}-\bm{s}),U(\bm{s}-\mu)\rangle^2}{\|U(\bm{r}-\bm{s})\|^2}
-\|U(\bm{s}-\mu)\|^2
\right).
\]
\end{block}

\begin{alertblock}{Key Decoupling Insight}
Modes of projected Gaussians are \textbf{independent of morphology} $U_n$.

This enables a sequential strategy:
\begin{enumerate}
\item estimate trajectories from projection modes,
\item estimate rotation and morphology from full sinogram fit.
\end{enumerate}
\end{alertblock}

\end{column}

\begin{column}{0.36\textwidth}

\begin{block}{Two-Stage Inference (GMM-CT)}
\textbf{Stage 1: Trajectory estimation}
\[
\Gamma_P^* = \argmin_{\{\eta_n\}_{n=1}^N}
\sum_{m_t=1}^{M_t}\sum_{n=1}^{N}
 d_H\bigl(M[g](t_{m_t}),\widehat{r}[\eta_n](t_{m_t})\bigr),
\]
where $M[g](t)$ are detected projection modes and $d_H$ is Hausdorff distance.

\vspace{0.4em}
Approximation details:
\begin{itemize}
\item Hungarian assignment (Jonker-Volgenant implementation),
\item multi-start L-BFGS,
\item Newton refinement for trajectory mode equations.
\end{itemize}

\vspace{0.5em}
\textbf{Stage 2: Rotation and morphology estimation}
\[
[\Gamma_R^*,\Gamma_S^*] =
\argmin_{\Gamma_R,\Gamma_S}
\frac{1}{M_tM_r}\sum_{\bm{m}\in\I}
H_\delta\!\left(g[\bm{m}] - F[\Gamma_P^*,\Gamma_R,\Gamma_S](\bm{m})\right),
\]
with Huber loss $H_\delta$.
\end{block}

\begin{block}{Algorithm Sketch}
\begin{center}
\begin{tikzpicture}[node distance=0.7cm]
\node[draw,rounded corners,fill=white,minimum width=13cm,minimum height=1.2cm] (in) {\textbf{Input:} dynamic sinogram, source-detector geometry, $N$};
\node[draw,rounded corners,fill=white,below=of in,minimum width=13cm,minimum height=1.2cm] (modes) {Detect sinogram peaks across time};
\node[draw,rounded corners,fill=white,below=of modes,minimum width=13cm,minimum height=1.2cm] (traj) {Optimize projectile parameters $\Gamma_P$ via mode matching};
\node[draw,rounded corners,fill=white,below=of traj,minimum width=13cm,minimum height=1.2cm] (init) {Initialize $\Theta$ (grid search) and $\alpha$ (NNLS)};
\node[draw,rounded corners,fill=white,below=of init,minimum width=13cm,minimum height=1.2cm] (joint) {Jointly optimize $(\Theta,\alpha,U)$ with Huber loss};
\node[draw,rounded corners,fill=white,below=of joint,minimum width=13cm,minimum height=1.2cm] (out) {\textbf{Output:} $\Gamma_P^*,\Gamma_R^*,\Gamma_S^*$};
\draw[-{Stealth[length=3mm]},thick] (in) -- (modes);
\draw[-{Stealth[length=3mm]},thick] (modes) -- (traj);
\draw[-{Stealth[length=3mm]},thick] (traj) -- (init);
\draw[-{Stealth[length=3mm]},thick] (init) -- (joint);
\draw[-{Stealth[length=3mm]},thick] (joint) -- (out);
\end{tikzpicture}
\end{center}
\end{block}

\begin{block}{Experimental Setup (2D)}
\begin{itemize}
\item $N=5$ Gaussian particles under gravity and independent rotation.
\item Fixed source at $(-1,-1)^\top$; 128 detector elements on a vertical line.
\item $M_t=150$ time samples on $[0,1.5]$ seconds.
\item Initialization: trajectory multi-starts ($N_{\mathrm{traj}}=20$),
rotation grid search ($N_{\mathrm{rot}}=200$),
morphology trials ($N_{\mathrm{morph}}=3$).
\end{itemize}
\end{block}

\end{column}

\begin{column}{0.32\textwidth}

\begin{block}{Numerical Results}
\centering
\includegraphics[width=0.92\linewidth]{\resdir/acquisition_geometry_exact.pdf}

\vspace{0.5em}
\includegraphics[width=0.92\linewidth]{\resdir/observed_sinogram.pdf}

\vspace{0.2em}
{\small Acquisition geometry and observed sinogram.}
\end{block}

\begin{block}{Mode-Based Trajectory Recovery}
\centering
\includegraphics[width=0.92\linewidth]{\resdir/projection_modes.pdf}

\vspace{0.5em}
\includegraphics[width=0.92\linewidth]{\resdir/trajectory_fitting_K5.png}

\vspace{0.2em}
{\small Detected modes and fitted projected trajectories.}
\end{block}

\begin{alertblock}{Reconstruction Quality}
\centering
\includegraphics[width=0.94\linewidth]{\resdir/individual_gaussian_reconstruction.pdf}

\vspace{0.4em}
For this case study, reconstructed particle morphology errors are on the order of $10^{-2}$ in absolute scale, with substantial improvement from initialization.
\end{alertblock}

\begin{block}{Takeaways and Next Steps}
\begin{itemize}
\item Joint motion-morphology inversion is feasible with closed-form GMM forward model.
\item Sequential decoupling improves robustness and optimization efficiency.
\item Framework extends naturally to $d=3$.
\item Ongoing work: noisy projections, unknown model order $N$, non-rigid dynamics.
\end{itemize}
\end{block}

\end{column}

\end{columns}
\end{frame}
\end{document}
